% KDD submission draft — compile on Overleaf with ACM "sigconf" / acmart template.
% Local build: ensure acmart.cls is on TEXINPUTS (Overleaf includes it by default).
\documentclass[sigconf,review]{acmart}

\settopmatter{printacmref=false, printccs=false, printfolios=true}

\usepackage{booktabs}
\usepackage[table]{xcolor}
\usepackage{graphicx}
\usepackage{adjustbox}
\usepackage{placeins}
\usepackage{microtype}
\usepackage{seqsplit}
\usepackage{amsmath}
\usepackage{amssymb}
\usepackage{enumitem}
\usepackage{algorithm}
\usepackage{algpseudocode}
\usepackage{tikz}
\usetikzlibrary{positioning, shapes.geometric, arrows.meta, calc, fit, backgrounds}

\AtBeginDocument{%
  \providecommand\BibTeX{{Bib\TeX}}%
}

% Breakable inline identifiers (avoids overfull boxes on long\_names in two columns).
\newcommand{\code}[1]{\texttt{\seqsplit{#1}}}

% Prefer draw.io export (PDF/PNG); fall back to TikZ .tex when exports are absent.
\newcommand{\includefig}[3][\linewidth]{%
  \IfFileExists{#2.pdf}{%
    \includegraphics[width=#1]{#2}%
  }{%
    \IfFileExists{#2.png}{%
      \includegraphics[width=#1]{#2}%
    }{%
      \input{#3}%
    }%
  }%
}

\title{Customer-Service LLM Runtime Harness:\\
Governed Takeover, Clarification-Aware Evidence, and Trace-Driven Improvement}

\author{Congde Yuan}
\email{congdeyuan@gmail.com}
\affiliation{%
  \city{Shenzhen}
  \country{China}
}

\author{Jun Wu}
\email{john.wuj@outlook.com}
\affiliation{%
  \city{Shenzhen}
  \country{China}
}

\author{Dafei Qiu}
\email{1643277704@qq.com}
\affiliation{%
  \city{Shenzhen}
  \country{China}
}

\author{Mingo Wu}
\email{mingo-wu@outlook.com}
\affiliation{%
  \city{Shenzhen}
  \country{China}
}

\author{Charles Quan}
\email{quanzhch@gmail.com}
\affiliation{%
  \city{Shenzhen}
  \country{China}
}

\author{Suipeng Li}
\email{suipeng_lee@foxmail.com}
\affiliation{%
  \city{Shenzhen}
  \country{China}
}

\author{Wu Mo}
\email{mowu_work@foxmail.com}
\affiliation{%
  \city{Shenzhen}
  \country{China}
}

\begin{document}

\begin{abstract}
Production customer-service bots must let large language models handle open-ended user language without giving them unbounded authority over account tools, policy rules, and human escalation. We present a \textbf{customer-service LLM runtime harness} deployed in a real-world workflow over deposit, withdrawal, P2P, KYC, asset, and account scenes. The system is implemented as a fixed LangGraph DAG~\cite{langgraph2024} with up to two structured-JSON LLM decisions on retrieval turns (classify, then elect), while the runtime enforces tool permissions, slot clarification, evidence assembly, rule governance, degradation, and audit logging. The paper contributes three reusable deployment patterns: \textbf{D1, governed runtime harness}, where the LLM owns semantic decisions but the runtime owns execution boundaries; \textbf{D2, clarification-aware evidence governance}, where FAQ hits, live account-tool outputs, and policy rules become typed, auditable evidence, and clarification is treated as a first-class bounded action; and \textbf{D3, runtime readiness and trace-driven improvement funnel}, where traces diagnose whether failures come from recall, ranking, candidate election, clarification, rules, tools, or action policy.
On $N{\approx}309$ labeled FAQ sessions (issue-ID Hit@$K$ only; Section~\ref{sec:eval-scope}), \textbf{canonical offline evidence} comes from Runs~1--3: stock rerank reaches $\approx$57\% Hit@1 while candidate election with top-20 candidates reaches $\approx$81\% ($<$1pp across Qwen3.5-27B vs.\ GPT-4o); a live top-10 ablation lowers this to 78.96\% (Table~\ref{tab:exp-topk}). Fine-tuning \code{bge-reranker-large}~\cite{xiao2024bge} lifts rerank P@1 to 92.9\% at $\approx$780\,ms mean latency; measured end-to-end joint under the fine-tuned checkpoint reaches \textbf{85.8\%} (Run~4), with top-10 and top-20 windows identical once rerank@10 saturates. A stock-rerank joint script ($N{=}303$; 67.99\% LLM Hit@1) remains a reproducibility reference in Section~\ref{sec:canonical-metrics}. Section~\ref{sec:deployment} specifies post-launch metrics from the production traffic window. The main takeaway: \emph{diagnose the runtime funnel and improve evidence quality before expanding LLM authority}---not backbone scaling alone.
\end{abstract}

\keywords{large language models, customer service, retrieval-augmented generation, hybrid retrieval, production systems}

\maketitle

\section{Introduction}
Customer-service assistants in regulated, high-stakes domains combine heterogeneous evidence sources: unstructured FAQ articles, structured backend tool outputs (order status, KYC state, asset balances), and deterministic business rules. Existing systems typically adopt one of two extremes:
\begin{itemize}[leftmargin=1.2em]
  \item \textbf{Pure RAG pipelines}~\cite{lewis2020rag,gao2024ragsurvey} retrieve semantically similar FAQ entries and generate answers with an LLM. They handle paraphrase and long-tail wording well, but may ignore account-specific facts and strict policy constraints when backend evidence is available.
  \item \textbf{Hard rule-first pipelines}~\cite{young2013pomdp,garcez2023neurosymbolic} invoke backend services and return rule-matched answers directly. They enforce policy precisely when triggers fire, but over-trigger on noisy tool outputs, cannot reconcile conflicting evidence, and fail closed when a channel times out.
\end{itemize}

We study a deployed LLM-based takeover of a customer-service workflow over deposit, withdrawal, P2P, KYC, asset, and account queries.
The central system object is a \emph{runtime harness}: the LLM is allowed to make semantic decisions, but serving-time code owns the boundaries around evidence, tool access, state, clarification, escalation, degradation, and audit.
Within the harness, FAQ retrieval results, backend account facts, and symbolic policy hits are normalized into typed evidence records with provenance, rank signals, and action constraints.
The LLM does not freely plan tool calls or synthesize unsupported answers; it classifies the scene, requests clarification when evidence is incomplete, and elects a discrete issue ID and dialogue action over a bounded evidence pool.
The implementation is a LangGraph~\cite{langgraph2024} \code{StateGraph} with six nodes (\code{analyze}, \code{guest\_mode}, \code{classify\_scene}, \code{kb\_only\_retrieval} or \code{parallel\_retrieval}, \code{llm\_result})---a fixed DAG, not an open-ended tool loop~\cite{yao2023react,singh2025agenticrag}.
When retrieval runs, the LLM is invoked twice: classify (\code{cs\_classify\_scene\_rules}) then elect (\code{cs\_llm\_result}).
Retrieval, MCP tools, and JSON rules supply evidence; Python enforces hard guardrails (guest isolation, slot caps, escalation); the LLM performs bounded election over fused evidence.

\paragraph{Reusable deployment patterns (D1--D3).}
\label{contrib:controller}
\label{contrib:ranking}
\textbf{D1: governed runtime harness.} Relative to retrieve-then-generate RAG~\cite{lewis2020rag} and single-loop tool agents~\cite{yao2023react}, we treat the LLM as a bounded decision component inside a production runtime, not as the owner of the whole action loop. The model emits structured JSON for scene classification and issue/action election; the runtime enforces tool eligibility, required-slot clarification, guest isolation, transfer caps, timeouts, fallbacks, and trace logging. This matches the emerging view that practical LLM agents are shaped by the harness around the model, not only by model weights~\cite{zhou2026externalization}. The reusable idea is that customer-service takeover requires a serving-time contract for what the model may see, ask, call, choose, and log.
\textbf{D2: clarification-aware evidence governance.} Inspired by ReAct's interleaving of reasoning and acting~\cite{yao2023react}, we treat clarification as a first-class runtime action alongside retrieval, tool calls, answer election, transfer, and blocking. Unlike a free-form ReAct loop, the action space is governed: missing MCP arguments trigger \textsc{Clarify}, prior \code{classify\_json} carries slots across turns, counters bound repeated clarification, and unresolved account-specific uncertainty escalates. FAQ rows, MCP outputs, and symbolic rule hits are normalized into typed evidence; rules become auditable \code{category=rule} evidence with \code{rule\_answer=None} rather than hard control-flow endpoints. This is a production-oriented evolution of reasoning--acting: the model can ask and decide, but runtime state decides when asking is still allowed.
\textbf{D3: runtime readiness and trace-driven improvement funnel.} LLM takeover quality is diagnosed as a runtime funnel rather than a single answer score: scene routing, evidence recall, rerank head quality, candidate election, clarification/action correctness, and online business outcomes. Runs~1--3 cover the evidence/ranking/election portion: stock rerank reaches $\approx$57\% Hit@1, candidate election reaches $\approx$81\%, changing the LLM backbone changes $<$1pp, and in-domain rerank fine-tuning reaches 92.9\% P@1; measured FT joint reaches 85.76\% (+4.5pp vs.\ Run~1; Table~\ref{tab:exp-topk}). The same traces support targeted updates to RAG queries, KB entries, reranker negatives, rule conditions, slot schemas, prompts, and transfer thresholds. We therefore do not use RAG as an autonomous self-evolution mechanism, but as one update target in a replayable data flywheel for deployed support systems~\cite{zhao2025agentloop}.

\paragraph{Design vs.\ engineering.}
Table~\ref{tab:design-vs-eng} separates reusable runtime patterns from implementation mechanisms.
We retain labels D1--D3 for cross-referencing: D1/D2 describe the deployed runtime harness and evidence-governance design; D3 describes the evaluation and improvement loop used to make the takeover measurable.

\paragraph{Contributions.}
Our contributions to applied customer-service LLM systems are:
\begin{itemize}[leftmargin=1.2em]
  \item \textbf{Governed runtime harness for LLM customer-service takeover (D1).} We report a deployed fixed-DAG harness in which the LLM owns semantic decisions while deterministic runtime code owns permissions, state, clarification bounds, escalation, degradation, and audit.
  \item \textbf{Clarification-aware evidence governance (D2).} We make clarification a governed action and compile FAQ, live account tools, and policy rules into typed auditable evidence for bounded candidate election; rule hits preserve symbolic provenance without brittle short-circuit returns.
  \item \textbf{Runtime readiness and trace-driven improvement funnel (D3).} We decompose takeover quality into stage-wise metrics and traces. On $N{\approx}309$ FAQ sessions, candidate election improves issue-ID Hit@1 from $\approx$57\% stock rerank to $\approx$81\%, domain rerank fine-tuning reaches 92.9\% P@1 without latency regression, and measured FT joint reaches 85.76\% (+4.5pp vs.\ Run~1); candidate-window width matters under stock rerank but not once rerank@10 saturates (Table~\ref{tab:exp-topk}). These traces define where RAG, rerank, rules, slots, prompts, and transfer policy should be updated after launch.
\end{itemize}

\begin{table}[t]
  \caption{Reusable runtime patterns vs.\ production mechanisms. D1--D2 are implemented in the deployed harness; D3 is partially measured offline today and extended with post-launch metrics in Section~\ref{sec:deployment}.}
  \label{tab:design-vs-eng}
  \centering
  \footnotesize
  \setlength{\tabcolsep}{2pt}
  \adjustbox{width=\columnwidth,center}{%
  \begin{tabular}{@{}>{\raggedright\arraybackslash}p{0.14\columnwidth}>{\raggedright\arraybackslash}p{0.40\columnwidth}>{\raggedright\arraybackslash}p{0.40\columnwidth}@{}}
    \toprule
    & \textbf{Reusable pattern} & \textbf{Production mechanisms} \\
    \midrule
    D1 &
      Governed runtime harness for LLM takeover
      & LangGraph DAG; dedicated \code{kcbot\_llm\_*}; MCP prompt hot-load; guest KB-only; \code{cs\_enabled\_tools}; slot \textsc{Clarify}; \code{\_chain\_log} \\
    D2 &
      Clarification-aware evidence governance
      & \code{parallel\_retrieval}; \textsc{GatherTimeout}; MCP aggregator; JSON rule matcher; \code{rule\_answer=None}; clarify counters \\
    D3 &
      Runtime readiness and trace-driven improvement funnel
      & Stage logs; Hit@$K$ funnel; rerank FT; planned rule/action labels; feature flags; replayable artifacts \\
    \bottomrule
  \end{tabular}}
\end{table}

The central claim is therefore \textbf{governed takeover, not autonomous takeover}: the LLM can own semantic judgment only because the runtime harness owns evidence, tools, clarification, escalation, and observability.
Figure~\ref{fig:classify_control_flow} illustrates the governed workflow; Related Work positions it against RAG-only, rule-first, and open agent loops---not a general autonomous-agent benchmark.

\paragraph{What we evaluate now vs.\ what the system supports.}
Section~\ref{sec:experiments} reports complementary offline studies on $N{\approx}309$ labeled FAQ sessions (issue-ID Hit@$K$)---sufficient to stress-test the \textbf{evidence/ranking/election slice} of the runtime funnel (D3) but \textbf{not} a full end-to-end score of clarification, MCP routing, rule triggers, action labels, or online business impact.
\textbf{Canonical metrics} for headline claims are Runs~1--3 (Table~\ref{tab:canonical-metrics}); Run~4 reports measured end-to-end joint accuracy under fine-tuned rerank plus candidate-window sensitivity (Sections~\ref{sec:canonical-metrics},~\ref{sec:exp-joint}).
Full runtime action metrics remain planned for future evaluation (Section~\ref{sec:exp-planned}).

\section{Related Work}
RAG grounds LLM answers in retrieved documents~\cite{lewis2020rag,gao2024ragsurvey}; adaptive variants decide when to retrieve or critique retrieved passages~\cite{asai2024selfrag,jeong2024adaptiverag,yan2024crag}. Customer-service RAG has also been deployed with structured knowledge, e.g., ticket knowledge graphs at LinkedIn~\cite{xu2024cskg}. These systems improve grounding but usually keep retrieval as the main control surface. Our system instead treats retrieval as one evidence channel inside a runtime that also governs account tools, policy rules, clarification, escalation, degradation, and trace logging.

Hybrid lexical--semantic retrieval and reranking are standard ways to improve candidate quality~\cite{robertson2009bm25,chen2022hybrid,bruch2023fusion,nogueira2019bert,xiao2024bge}. We use these components pragmatically: three-channel FAQ recall supplies a broad pool, optional rerank sharpens head precision, and production fallbacks return RRF order when rerank fails. The key applied lesson is not that reranking exists, but that the runtime funnel reveals when rerank quality is the first bottleneck: on our labeled sessions, in-domain rerank adaptation improves P@1 more than switching the decision LLM.

Tool-augmented agents such as ReAct, Toolformer, and ToolLLM let LLMs interleave reasoning with API calls~\cite{yao2023react,schick2023toolformer,qin2023toolllm}. Enterprise customer service needs a narrower contract: tools are scene-specific, account-bound, and policy-sensitive. Related constrained systems include DFA-RAG, which uses finite-state routing for regulated dialogue~\cite{sun2024dfarag}, and task-oriented dialogue work, which separates state, slots, and action policy~\cite{young2013pomdp,wen2017network,chung2023instructtods}. Our D2 contribution can be read as a production evolution of ReAct: the model may ask or act, but slot sufficiency, repeated clarification, tool eligibility, and human transfer are governed by runtime state.

Finally, recent work argues that LLM systems improve through externalized memory, protocols, and harnesses rather than model weights alone~\cite{shinn2023reflexion,zhou2026externalization}. Customer-support Agent-in-the-Loop shows that operational feedback can form a data flywheel~\cite{zhao2025agentloop}. Our D3 follows this direction, but keeps the loop auditable: traces create review queues, reranker negatives, KB backlogs, rule-test fixtures, and replay sets rather than directly rewriting production behavior at serving time.

\section{Problem Setting}
Given a user request $q$, conversation history $H$, and optional user state $u$ (authentication, site locale), the system predicts a response action $a$ and response content $y$.

Actions follow task-oriented dialogue conventions~\cite{young2013pomdp,wen2017network}:
\begin{itemize}[leftmargin=1.2em]
  \item \textsc{Answer}: return a grounded response from a selected FAQ issue,
  \item \textsc{Clarify}: ask for missing slots or intent clarification,
  \item \textsc{Human\_Transfer}: escalate to a human agent,
  \item \textsc{Chitchat} / \textsc{Blocked}: non-business or policy-blocked handling.
\end{itemize}

The system must jointly satisfy:
\begin{itemize}[leftmargin=1.2em]
  \item \textbf{Factual correctness} grounded in KB or backend evidence,
  \item \textbf{Business-policy compliance} via rule grounding without over-triggering,
  \item \textbf{Multi-turn robustness} under incomplete slot filling,
  \item \textbf{Graceful degradation} when one retrieval or tool channel fails.
\end{itemize}

\section{Method}
\subsection{Notation and Objectives}
\label{sec:notation}
Let $q^{(t)}$ denote the latest user utterance at turn $t$, $H$ the trimmed dialogue history, and $u$ user state (login flag, locale). The classify-stage LLM outputs a structured tuple
\begin{equation}
\label{eq:classify}
  \Gamma(q^{(t)}, H) = \bigl(t,\; b_{\text{tool}},\; \mathbf{s},\; \tilde{q},\; a_{\text{cls}}\bigr),
\end{equation}
where $t \in \mathcal{T}$ is a scene/tool label ($|\mathcal{T}|{=}7$ in production), $b_{\text{tool}} \in \{0,1\}$ is \code{need\_tool\_call}, $\mathbf{s}$ slot values for MCP arguments, $\tilde{q}$ a rewritten retrieval query, and $a_{\text{cls}}$ an optional early action. The decision-stage LLM maps a candidate pool $\mathcal{C}$ to
\begin{equation}
\label{eq:decide}
  \Delta(\mathcal{C}, q^{(t)}, H) = \bigl(\hat{\iota},\; \hat{a},\; \hat{y}\bigr),
\end{equation}
with $\hat{\iota}$ an elected FAQ \code{issue\_id}, $\hat{a}$ a dialogue action, and $\hat{y}$ response text. Our goal is to maximize grounded issue/action accuracy while preserving hard guardrails (guest isolation, slot clarification, escalation) independent of $\Delta$.

\subsection{LLM-Centric Pipeline Design}
\label{sec:llm-centric}
\noindent\textit{Implements D1 runtime harness (classify stage).}

The framework follows an \emph{LLM-in-the-loop} design implemented as a LangGraph~\cite{langgraph2024} \code{StateGraph}. Every non-terminal turn invokes the LLM twice when retrieval runs. The six graph nodes are:
\begin{center}
  \footnotesize
  \code{analyze} $\rightarrow$ \code{guest\_mode} $\rightarrow$ \code{classify\_scene} $\rightarrow$\\
  \{\code{kb\_only\_retrieval} $|$ \code{parallel\_retrieval} $|$ early exit\} $\rightarrow$ \code{llm\_result}
\end{center}
\begin{itemize}[leftmargin=1.2em]
  \item \textbf{Classify-stage LLM} (\code{llm\_classify\_scene}; prompt \code{cs\_classify\_scene\_rules}): scene/tool selection over seven labels, slot extraction, \code{need\_tool\_call} routing, and \code{rewritten\_query}. Early actions include \textsc{Chitchat}, \textsc{Blocked}, and \textsc{Human\_Transfer}. JSON is parsed via \code{response\_format} with repair on malformed output. Multimodal images use \code{image\_url} blocks when present.
  \item \textbf{Decision-stage LLM} (\code{node\_llm\_result}; prompt \code{cs\_llm\_result}): issue/action election over KB and rule candidates in \code{<kb\_entry>} and \code{<account\_rule\_match>} tags, yielding \code{issue\_id}, \code{summary}, and \code{reason}. Optional \code{cs\_answer\_summary\_details} adds a formatted answer with collapsible KB source.
\end{itemize}
Both stages use a dedicated CS LLM endpoint (\code{kcbot\_openai\_api\_*}) isolated from the general copilot model.
Prompts are fetched from an MCP prompt server with Redis cache and local file fallback.
Non-LLM modules---hybrid retrieval, MCP tool calls, JSON rule matching---execute \emph{between} these two stages and supply evidence rather than dictating the final outcome.

The production choice is deliberately not ``give the model all tools and ask it to solve the case.'' A customer-service turn has several authorities that must remain separable: identity authority (guest vs.\ logged-in), tool authority (which backend can be queried), evidence authority (which facts can ground the answer), action authority (answer, clarify, transfer, block), and operational authority (timeouts, fallback, logging). The runtime harness decomposes these authorities into deterministic gates around two LLM decisions. This is the core reusable unit: another support domain can replace scenes, tools, rules, and FAQ schema while preserving the same governed takeover pattern.

\subsection{Workflow Overview}
\label{sec:workflow}
\noindent\textit{Production workflow (Table~\ref{tab:design-vs-eng}, D1).}
The \code{analyze} node trims dialogue history (at most 21 turns and 20{,}000 characters) and maps platform locale to \code{reply\_language}.
\code{guest\_mode} sets \code{is\_guest\_mode} from \code{user\_id}.
\code{classify\_scene} may short-circuit to \textsc{END} on special intents or slot \textsc{Clarify}.
Otherwise \code{route\_after\_classify} sends \emph{guests} to \code{kb\_only\_retrieval}; all \emph{authenticated} users enter \code{parallel\_retrieval} (FAQ recall always runs; MCP/rule matching runs only when \code{need\_tool\_call=true} inside that node).
Both retrieval paths converge on \code{llm\_result}.

\begin{figure*}[t]
  \centering
  \adjustbox{max width=0.86\textwidth,max height=0.52\textheight,center}{%
    \includefig[\textwidth]{figures/fig1_classify_control_flow}{figures/fig1_classify_control_flow}%
  }
  \caption{Classify-stage control flow: early actions, MCP slot handling, and guest/member routing to \texttt{llm\_result}.}
  \label{fig:classify_control_flow}
\end{figure*}

Figure~\ref{fig:classify_control_flow} expands classify-stage branching, post-classify routing, and decision-stage outcomes. The DAG enables per-node timeouts, logging, and partial-result preservation.

\subsection{LLM-Conditioned Dynamic Routing}
\label{sec:routing}
\noindent\textit{Implements D1 runtime guardrails (Table~\ref{tab:design-vs-eng}).}

Equation~\eqref{eq:classify} is implemented by \code{llm\_classify\_scene}. Let $\mathcal{A}_{\text{early}}=\{\textsc{Chitchat},\textsc{Blocked},\textsc{Human\_Transfer}\}$. Graph routing in \code{route\_after\_classify} is a deterministic function of guest mode and $a_{\text{cls}}$:
\begin{equation}
\label{eq:route}
  \text{next}(u, a_{\text{cls}}) =
  \begin{cases}
    \textsc{End}, & a_{\text{cls}} \in \mathcal{A}_{\text{early}} \lor \text{slot \textsc{Clarify}},\\[2pt]
    \text{\code{kb\_only\_retrieval}}, & \text{guest}(u) = 1,\\[2pt]
    \text{\code{parallel\_retrieval}}, & \text{logged-in}(u).
  \end{cases}
\end{equation}
Authenticated members therefore always enter \code{parallel\_retrieval}; $b_{\text{tool}}$ gates the \emph{structured} branch inside that node rather than the graph edge.
Let $d_{\text{mcp}} = \textsc{MCP}(t, \mathbf{s})$ when $b_{\text{tool}}{=}1$ and the user is logged in; otherwise the MCP branch is skipped (\code{skip\_mcp\_no\_tool\_call} or \code{skip\_mcp\_no\_user\_id}) while FAQ recall still runs in parallel with an empty tool payload.
Programmatic guardrails complement LLM outputs.
If the LLM selects a tool disabled by \code{cs\_enabled\_tools}, the scene falls back to \code{other} with \code{need\_tool\_call=false}.
When MCP marks arguments as required but slots are missing, Python returns slot \textsc{Clarify} (up to \code{slot\_clarify\_count}${\leq}3$) \emph{before} retrieval.
Prior-turn \code{classify\_json} in \code{metaData} is fed back as assistant content so slots accumulate across turns; switching tools resets \code{slot\_clarify\_count} and accumulated arguments to avoid cross-scene contamination.

\subsection{Dual-Path Retrieval with Rule-as-Evidence Injection}
\label{sec:rule-fusion}
\noindent\textit{Implements D2 (clarification-aware evidence governance; Section~\ref{sec:rule-fusion}).}

For \code{parallel\_retrieval}, FAQ and MCP tasks are launched concurrently; MCP tools are exposed via the Model Context Protocol (MCP)~\cite{modelcontextprotocol2024} and aggregated by a scene-specific client:
\begin{equation}
\label{eq:parallel}
  (\mathcal{C}_{\text{faq}}, d_{\text{mcp}}) = \textsc{GatherTimeout}\!\left(
    \textsc{HybridRetrieve}(\tilde{q}),\;
    \mathbb{1}[b_{\text{tool}} \wedge \text{login}(u)] \cdot \textsc{MCP}(t,\mathbf{s})
  \right),
\end{equation}
where \textsc{GatherTimeout} preserves whichever branch finishes within \code{cs\_parallel\_total\_timeout} (default 120s) and does not discard partial results on single-branch failure.

\paragraph{Illustrative rule-as-evidence case (D2).}
Consider a logged-in user asking whether a withdrawal is delayed while MCP returns \code{status=pending} and \code{delayReason=chain\_congestion}.
A rule-first bot might immediately return a canned ``wait for chain confirmation'' answer (short-circuit).
In our pipeline, \textsc{Match} fires the corresponding policy rule, which is injected as a \code{category=rule} candidate \emph{ahead of} FAQ hits about generic withdrawal FAQs.
The decision-stage LLM still reads the latest user wording and may elect a more specific \code{issue\_id} (e.g., ``accelerate review'' vs.\ ``explain congestion'') or return \textsc{Clarify} if slots are incomplete---preserving auditability of symbolic triggers without giving up semantic disambiguation~\cite{garcez2023neurosymbolic,young2013pomdp}.
We do not yet report offline precision/recall for rule triggers.

\paragraph{Symbolic rule matching.}
Given normalized MCP JSON $\phi(d_{\text{mcp}})$, each policy rule $r$ defines conjunctive conditions (JSON \code{\$and} arrays with nested \code{\$or}/\code{\$expr}). We write
\begin{equation}
\label{eq:rule-match}
  \textsc{Match}(r,\phi) = \bigwedge_{k} \textsc{EvalSub}_k\bigl(\phi\bigr),
\end{equation}
where \textsc{EvalSub} implements field resolvers and operators (\code{contains}, numeric compares, \code{\_fallback}) in \code{customer\_service\_rule\_matcher}. With \code{return\_all=true},
\begin{equation}
\label{eq:rule-set}
  \mathcal{R} = \{ r : \textsc{Match}(r,\phi) \}.
\end{equation}
Each $r \in \mathcal{R}$ is converted to a KB-shaped candidate record ( \code{category=rule} ) and prepended:
\begin{equation}
\label{eq:candidate-merge}
  \mathcal{C} = \textsc{RulesToCandidates}(\mathcal{R}) \,\|\, \mathcal{C}_{\text{faq}},
\end{equation}
where $\|$ denotes list concatenation with rule rows first. Production always sets \code{rule\_answer=None}, so even $|\mathcal{R}|{=}1$ flows to $\Delta$ in~\eqref{eq:decide} rather than a hard short-circuit. Escalation without election applies when MCP data exist, $\mathcal{R}=\emptyset$, and prior slot clarification count exceeds zero (Section~\ref{sec:decision}).

\paragraph{Unified candidate contract.}
The novelty of D1/D2 is not the existence of multiple retrievers, but the runtime contract that all evidence must satisfy before the decision LLM can use it.
Each candidate in $\mathcal{C}$ carries: (i) a stable \code{issue\_id} or rule-derived pseudo-ID, (ii) a \code{category} such as FAQ or rule, (iii) provenance fields identifying the retrieval channel, rule ID, or MCP scene, (iv) rank signals such as RRF score, rerank score, and candidate position when available, and (v) an allowed action surface (\textsc{Answer}, \textsc{Clarify}, \textsc{Human\_Transfer}) rather than an unconstrained response.
This contract has three practical effects.
First, it makes heterogeneous evidence comparable without forcing backend JSON into a free-form prompt.
Second, it supports graceful degradation: if dense retrieval, rerank, or MCP fails, the remaining candidates still satisfy the same interface.
Third, it makes the decision trace auditable because the final \code{issue\_id} and action can be tied back to an explicit candidate record.
The decision prompt therefore asks the LLM to resolve ambiguity among grounded candidates, not to discover hidden tools or infer business policy from raw logs.

\paragraph{Clarification as governed ReAct.}
Clarification is handled as an action over state, not as arbitrary model initiative. If the classify-stage LLM identifies a scene but required MCP slots are absent, the runtime asks a slot-specific question before any account-bound tool call. If the decision-stage LLM cannot choose among grounded candidates, it may emit \textsc{Clarify}, but \code{clarify\_count} bounds repeated turns and can force \textsc{Human\_Transfer}. This keeps the useful ReAct pattern---reason about missing information, then act by asking or calling a tool~\cite{yao2023react}---while making the ask/call boundary auditable and reversible in logs.

Rule-as-evidence injection is therefore an evidence-construction mechanism; the decision-stage LLM performs grounded election over the unified pool.

\subsection{Hybrid Retrieval and Ranking}
\label{sec:hybrid}
\noindent\textit{Implements D3 runtime funnel (hybrid recall and rerank slice).}

FAQ retrieval in \code{hybrid\_recall.py} executes three channels, with BM25 and embedding fetches started in parallel:
\begin{itemize}[leftmargin=1.2em]
  \item BM25 over \code{issue}$^{2}$ and \code{issueDesc} (fetch 150),
  \item KNN on \code{text\_vector} (issue embedding, fetch 10),
  \item KNN on \code{desc\_vector} (issueDesc embedding, fetch 10).
\end{itemize}
Failed channels are logged and skipped; remaining hits feed \code{merge\_rrf\_hits}.
For channel $c \in \{\text{bm25}, \text{text}, \text{desc}\}$, let $L_c = (d_{c,1}, d_{c,2}, \ldots)$ be the ranked hit list returned by OpenSearch.
BM25 uses weighted \code{multi\_match} on \code{issue}$^{2}$ and \code{issueDesc}; KNN channels query $\mathbf{e}(\tilde{q})$ against \code{text\_vector} and \code{desc\_vector}.

Following \code{merge\_rrf\_hits}, each document $d$ (keyed by \code{issue\_id}) receives a weighted reciprocal rank fusion score with per-channel deduplication (a channel contributes at most once per ID):
\begin{equation}
\label{eq:rrf}
  S_{\text{RRF}}(d) = \sum_{c \,:\, d \in L_c} \frac{w_c}{K + \mathrm{rank}_c(d)},
\end{equation}
with production defaults $K{=}40$, $w_{\text{bm25}}{=}1.05$, $w_{\text{text}}{=}0.85$, $w_{\text{desc}}{=}1.0$ from \code{RRF\_RANK\_CONSTANT} and \code{RRF\_W\_*} in code~\cite{cormack2009rrf,chen2022hybrid,bruch2023fusion}. Candidates are sorted by $S_{\text{RRF}}$ and truncated to $N_{\text{rrf}}{=}50$ (\code{CS\_KB\_RRF\_TOP\_N}).

\paragraph{Cross-encoder rerank.}
Let $\pi_{\text{rrf}}$ be the RRF ordering. An optional rerank API scores pairs $(\tilde{q}, \text{doc}(d))$ with cross-encoder $s_{\text{ce}}$ and induces
\begin{equation}
\label{eq:rerank}
  \pi_{\text{out}} =
  \begin{cases}
    \mathrm{argsort}_{d}\, s_{\text{ce}}\bigl(\tilde{q}, \text{doc}(d)\bigr) & \text{rerank OK},\\[2pt]
    \pi_{\text{rrf}}[:k] & \text{timeout / error},
  \end{cases}
\end{equation}
matching \code{\_async\_rerank\_hits\_api} fallback behavior.
Final FAQ hits are $\mathcal{C}_{\text{faq}} = \pi_{\text{out}}[:k]$.
Offline evaluation (Run~3, Section~\ref{sec:exp-rerank}) shows that stock \code{bge-reranker-large}~\cite{xiao2024bge,reimers2019sentencebert} achieves only 57.3\% business P@1 after recall, whereas \code{out\_bge\_rerank\_ft} reaches 92.9\% P@1 at $\approx$780\,ms mean latency---motivating the fine-tuned checkpoint in production.

Multilingual support is handled by language-specific \code{reply\_*} and \code{issue\_*} fields in the index, selected at query time based on detected user language.

\FloatBarrier
\subsection{State-Constrained LLM Decision Layer}
\label{sec:decision}
\noindent\textit{Implements D1 runtime harness (decision stage).}

The decision stage implements~\eqref{eq:decide} via \code{node\_llm\_result}: $(\hat{\iota}, \hat{a}, \hat{y}) = \Delta(\mathcal{C}, q^{(t)}, H)$ where $\mathcal{C}$ follows~\eqref{eq:candidate-merge} on dual-path turns or $\mathcal{C}=\mathcal{C}_{\text{faq}}$ on guest KB-only turns. Runs~1--2 (Section~\ref{sec:exp-results}) evaluate this mapping in isolation: with stock rerank, Hit@1 rises from $\approx$57\% to $\approx$81\% when the LLM reads tagged \code{<kb\_entry>} / \code{<account\_rule\_match>} blocks---evidence that $\Delta$ should not be replaced by $\pi_{\text{out}}[1]$ alone.

\paragraph{Programmatic state and escalation.}
Turn metadata $\mathcal{M}_t$ stored in \code{cs\_turn\_meta} augments the LLM with counters persisted to \code{metaData}:
\begin{equation}
\label{eq:state}
  \mathcal{M}_{t+1} = \mathcal{F}\bigl(\mathcal{M}_t,\; a_{\text{cls}},\; \hat{a},\; t\bigr),
\end{equation}
where $\mathcal{F}$ increments \code{slot\_clarify\_count} (MCP argument loops, cap 3 before forced transfer), \code{clarify\_count} (business clarification from $\hat{a}$), and merges \code{accumulated\_arguments}; tool switches reset slot counters to avoid cross-scene leakage. Hard transfer triggers include
\begin{equation}
\label{eq:transfer}
\begin{aligned}
  \textsc{Human\_Transfer} \;\Leftarrow\;&
    \bigl(\code{clarify\_count} \geq 3\bigr) \vee \bigl(\code{transfer\_req} \geq 2\bigr) \\
  &\vee \bigl(\text{MCP hit} \wedge \mathcal{R}=\emptyset \wedge \code{slot\_clarify\_count}>0\bigr).
\end{aligned}
\end{equation}
The \code{analyze} node enforces $|H'| \leq 21$ turns and $|\text{chars}(H')| \leq 20{,}000$ before any LLM call.
On \textsc{Answer}, the system resolves reply text from the elected \code{issue\_id}, optionally wraps it with LLM summary + \code{<details>} KB source, and may persist QA pairs when \code{cs\_qa\_persist\_enable} is on.

\section{System Deployment}
\label{sec:deployment}
The harness is deployed as a non-streaming API integrated with platform identity, risk controls, and observability infrastructure. The deployment is not a demo wrapper around an LLM: every production turn passes through runtime gates for user mode, tool eligibility, slot sufficiency, evidence construction, action selection, degradation, and logging. We emphasize deployment lessons that complement the offline FAQ metrics:

\paragraph{Channel isolation and timeouts.}
MCP tool calls and OpenSearch retrieval run under separate per-branch and global timeouts. If one branch fails, the other branch's results are preserved rather than discarding the entire request.

\paragraph{Configuration-driven degradation.}
Feature flags control OpenSearch recall, legacy KB API fallback, rerank enablement, and answer formatting modes. Partial outages degrade to RRF-only ranking~\cite{cormack2009rrf} or KB-only routing without full service interruption. The production rerank slot uses \code{out\_bge\_rerank\_ft} (Run~3) rather than the stock \code{bge-reranker-large} checkpoint~\cite{xiao2024bge}, reflecting the measured +35.6pp P@1 gain on labeled business traffic at no latency penalty. Offline Run~4 reports measured end-to-end accuracy under fine-tuned rerank (Table~\ref{tab:canonical-metrics}).

\paragraph{Auditability.}
Each request emits a chain log tracing classification JSON, MCP status, rule hit list, FAQ recall scores, and final decision rationale---supporting post-hoc review in regulated customer-service settings.

\paragraph{Trace-driven improvement loop.}
Runtime traces are designed to localize failures to a specific layer rather than treating wrong answers as generic LLM mistakes. A failed session can be attributed to scene routing, missing-slot clarification, tool availability, rule over-triggering, recall miss, rerank head error, candidate-election error, or action-policy error. This attribution determines the update target: KB content and query rewriting for recall misses, hard negatives or model choice for rerank errors, rule edits for policy mismatches, slot-schema or prompt edits for clarification failures, and transfer-threshold changes for action-policy failures. Run~3 is one concrete instance of this loop: logs exposed rerank head quality as the dominant bottleneck, leading to in-domain rerank adaptation rather than backbone scaling.

The loop is intentionally conservative. Production traces do not directly rewrite rules, prompts, or KB entries at serving time. Instead, they produce review queues and replay sets: ambiguous transfers become action-policy examples, missing knowledge becomes KB backlog, wrong FAQ elections become reranker negatives or prompt regression cases, and noisy rule fires become rule-test fixtures. This is closer to a customer-support data flywheel~\cite{zhao2025agentloop} than to autonomous self-modification. The resulting trade-off is slower than fully automatic self-evolution, but it preserves auditability, rollback, and controlled post-launch measurement.

\paragraph{Safety gates.}
Input and output guards, URL risk filtering, and image moderation run at API boundaries independent of the retrieval pipeline.

\paragraph{Post-launch measurement.}
Because the system is deployed, post-launch impact should be reported from production logs over a fixed traffic window and compared against the previous production configuration or a matched traffic slice. Table~\ref{tab:post-launch-placeholder} lists the deployment metrics to fill before submission; the values should quantify both business effect and runtime cost, and should state how online serving differs from the offline FAQ-only studies in Section~\ref{sec:experiments}.

\begin{table}[t]
  \caption{Post-launch deployment metrics. ``Baseline'' should be the pre-harness bot or the immediately preceding production configuration on the same traffic slice.}
  \label{tab:post-launch-placeholder}
  \centering
  \footnotesize
  \setlength{\tabcolsep}{3pt}
  \adjustbox{width=\columnwidth,center}{%
  \begin{tabular}{@{}>{\raggedright\arraybackslash}p{0.26\columnwidth}>{\raggedright\arraybackslash}p{0.24\columnwidth}>{\raggedright\arraybackslash}p{0.22\columnwidth}>{\raggedright\arraybackslash}p{0.20\columnwidth}@{}}
    \toprule
    \textbf{Metric} & \textbf{Baseline} & \textbf{Harness} & \textbf{Window} \\
    \midrule
    Human-transfer rate & TODO & TODO & TODO \\
    First-contact resolution proxy & TODO & TODO & TODO \\
    User re-contact / reopen rate & TODO & TODO & TODO \\
    End-to-end latency p50 / p95 & TODO & TODO & TODO \\
    Cost per resolved session & TODO & TODO & TODO \\
    Safety or policy escalation rate & TODO & TODO & TODO \\
    \bottomrule
  \end{tabular}}
\end{table}

\FloatBarrier
\section{Experiments}
\label{sec:experiments}

\subsection{Evaluation Scope and Claims}
\label{sec:eval-scope}

This paper is an \textbf{Applied Data Science} account of a deployed runtime harness: the method sections describe the full DAG (including MCP and rules), while the reported offline metrics intentionally isolate subsystems that can be labeled today. Runs~1--3 measure FAQ recall/rerank/election; clarification precision, rule-trigger precision, tool-call F1, action accuracy, and online SLOs are deployment-supported but not yet labeled in the offline set.

\paragraph{Why $N{\approx}309$ is still informative.}
The set is small relative to web-scale benchmarks but matches the \emph{labeled issue-ID inventory} available for supervised rerank adaptation and decision-stage audit in this domain---similar in spirit to curated CS evaluation sets in prior deployed work~\cite{xu2024cskg}. Runs are structured as \emph{controlled ablations} (vary one stage at a time) rather than a single leaderboard score, which partially offsets sample size. We report a 95/5 rerank split variant (89.3\% P@1) alongside full-data fine-tuning to surface optimistic bias (Section~\ref{sec:exp-rerank}).

Planned work extends to action-level and MCP/rule metrics (Section~\ref{sec:exp-planned}).

\subsection{Canonical Metrics and Run~4 Reproduction}
\label{sec:canonical-metrics}

Table~\ref{tab:canonical-metrics} distinguishes \textbf{which numbers to cite for which claim}. Runs~1--2 and Run~3 are logged from historical offline passes on $N{\approx}309$ cases and remain our \emph{primary} evidence for (i)~candidate-election lift under stock \code{bge-reranker-large} and (ii)~domain rerank fine-tuning.
Run~4 is a \textbf{measured end-to-end joint pass} under \code{out\_bge\_rerank\_ft} on $N{=}309$ cases (frozen CUDA rerank scores feeding \code{cs\_llm\_result}; artifact \code{Paper/data/run4\_joint\_ft\_measured.json}).
A separate stock-rerank joint script on $N{=}303$ CONTINUE-labeled cases (\code{paper\_faq\_pipeline\_eval.py}; 68.0\% LLM Hit@1) remains a reproducibility reference---lower than the Run~1 live archive (81.23\%) for harness and filter reasons discussed below---but is not the headline joint number.

\begin{table}[t]
  \caption{Canonical offline metrics: primary vs.\ replication vs.\ sensitivity.}
  \label{tab:canonical-metrics}
  \centering
  \footnotesize
  \setlength{\tabcolsep}{2pt}
  \adjustbox{width=\columnwidth,center}{%
  \begin{tabular}{@{}>{\raggedright\arraybackslash}p{0.12\columnwidth}>{\raggedright\arraybackslash}p{0.38\columnwidth}>{\raggedright\arraybackslash}p{0.44\columnwidth}@{}}
    \toprule
    \textbf{Source} & \textbf{Use in paper} & \textbf{LLM Hit@1} \\
    \midrule
    Runs~1--2 &
      Primary: election gain ($\approx$81\% vs.\ $\approx$57\% rerank) &
      \textbf{81.2\%} ($N{=}309$/308; top-20) \\
    Run~1 ablation &
      Candidate window under stock rerank (live repro) &
      \textbf{78.96\%} (top-10; Table~\ref{tab:exp-topk}) \\
    Run~4 (meas.) &
      Primary: FT joint end-to-end ($N{=}309$) &
      \textbf{85.8\%} (top-10/20; Table~\ref{tab:exp-topk}) \\
    Run~4 (stock ref.) &
      Replication: stock joint script ($N{=}303$) &
      68.0\% (not headline) \\
    \bottomrule
  \end{tabular}}
\end{table}

\paragraph{Why stock joint $\neq$ Run~1 live.}
Besides $N{=}303$ vs.\ $309$, the stock joint script (\code{paper\_faq\_pipeline\_eval.py}) differs from the historical Run~1 \code{workflow\_simple} live path in gateway routing, filters (\code{action=CONTINUE}), and stack revisions (Section~\ref{sec:exp-setup}).
We report both without retroactive relabeling of Run~1 archives.
Production deploys fine-tuned rerank (Run~3); Run~4 FT joint (85.76\%) is the measured complement to Run~1 stock election (81.23\%).

\subsection{Shared Dataset and Protocol}
\label{sec:exp-setup}

\textbf{Test sets.} We run two offline passes on labeled FAQ cases with gold \code{issue\_id} and multi-turn user history from production-style customer-service logs: Run~1 ($N{=}309$ valid cases) and Run~2 ($N{=}308$ valid cases). Case counts differ slightly due to filtering invalid labels between passes; the protocol is otherwise identical.

\textbf{Query construction.} User utterances are machine-translated to English. For \emph{retrieval}, all user turns in a session are concatenated into a single query string (chronological order) and passed to three-channel hybrid recall, which returns up to 50 FAQ candidates. For \emph{decision}, the decision-stage LLM (\code{cs\_llm\_result}) receives only the \emph{last} user turn as the active query, together with the recalled candidate list---mirroring production behavior where the latest utterance disambiguates intent while prior turns inform recall breadth.

\textbf{Decision models (Runs~1--2).} Run~1 uses Qwen3.5-27B; Run~2 uses Azure GPT-4o. Both emit structured JSON (\code{issue\_id}, action, rationale). Reranking uses the off-the-shelf \code{bge-reranker-large} checkpoint and reorders the recall top-50 before candidate election over \textbf{top-20} candidates (top-10 ablation in Table~\ref{tab:exp-topk}).

\textbf{Rerank protocol (Run~3).} User queries are translated to English and concatenated across turns; hybrid recall returns top-50 candidates, which are rescored by a cross-encoder reranker. We compare the public \code{bge-reranker-large} baseline against three fine-tuned variants trained on the same 309 issue-ID business pairs (query, positive FAQ, negatives from the recall pool): \code{out\_bge\_rerank\_ft} (all 309 examples), \code{out\_bge\_rerank\_ft\_split} (95\%/5\% train/eval split), and \code{out\_bge\_rerank\_ft\_mixed\_bizfull\_marco\_r20} (309 business pairs plus 2{,}000 randomly sampled MMarcoReranking training pairs~\cite{thakur2021beir}).

\textbf{Metric.} We report \textbf{Hit@$K$}: the fraction of cases whose gold \code{issue\_id} appears within the top-$K$ ranked candidates at a given pipeline stage. Hit@1 at the LLM stage is \emph{end-to-end issue-ID accuracy}---the model's elected \code{issue\_id} matches the label. At recall, Hit@50 indicates whether the gold issue appears anywhere in the retrieved pool (recall ceiling before rerank/LLM).

\subsection{LLM Decision Stage (Runs~1--2)}
\label{sec:exp-results}

\noindent\textit{Purpose.} Runs~1--2 test the election segment of the runtime readiness funnel (D3): retrieval and off-the-shelf rerank are treated as a \emph{candidate provider}, and the decision-stage LLM (\code{llm\_result}) performs structured issue-ID/action election inside the governed harness. We deliberately hold recall and rerank fixed (generic \code{bge-reranker-large}) and vary only the decision backbone to separate \emph{runtime/evidence gains} from \emph{model-scale gains}.

Table~\ref{tab:exp-summary} compares Hit@1 across both runs. Retrieval and rerank metrics are stable ($\pm$0.5pp), while decision-stage Hit@1 remains near 81\% for both LLMs---the $<$1pp gap between Qwen3.5-27B and GPT-4o implies that raising end-to-end accuracy on this benchmark requires better candidates (recall/rerank), not a larger generator alone. Table~\ref{tab:exp-hitk} reports full Hit@$K$ curves; Table~\ref{tab:exp-funnel} contrasts the stock pipeline measured in Runs~1--2 with the rerank stage after Run~3 fine-tuning (LLM omitted in Run~3). Figure~\ref{fig:performance_funnel} visualizes the Run~1 three-stage Hit@1 funnel.

\begin{figure}[t]
  \centering
  \includefig[\columnwidth]{figures/fig6_performance_funnel}{figures/fig6_performance_funnel}
  \caption{Run~1 issue-ID Hit@1 across pipeline stages ($N{=}309$; stock \code{bge-reranker-large} rerank; Qwen3.5-27B decision). Bars match Table~\ref{tab:exp-hitk} (Run~1 rows).}
  \label{fig:performance_funnel}
\end{figure}

\begin{table}[t]
  \caption{Issue-ID Hit@1 funnel across runs. Run~4: measured FT joint ($N{=}309$; Section~\ref{sec:exp-joint}). Run~3 rerank uses fine-tuned checkpoint on $N{\approx}309$.}
  \label{tab:exp-funnel}
  \centering
  \footnotesize
  \setlength{\tabcolsep}{2pt}
  \adjustbox{width=\columnwidth,center}{%
  \begin{tabular}{@{}lccc@{}}
    \toprule
    \textbf{Stage} & \textbf{Runs 1--2} & \textbf{Run 3} & \textbf{Run 4} \\
    \midrule
    Hybrid recall Hit@1   & 42.3 & 42.3 & 42.1 \\
    After rerank Hit@1    & 56.6 & 92.9 & 92.9 \\
    After decision Hit@1  & 81.5 & --- & \textbf{85.8} \\
    \bottomrule
  \end{tabular}}
\end{table}

\begin{table}[t]
  \caption{Issue-ID Hit@1 summary on labeled FAQ cases (same protocol per run).}
  \label{tab:exp-summary}
  \centering
  \footnotesize
  \setlength{\tabcolsep}{2.5pt}
  \adjustbox{width=\columnwidth,center}{%
  \begin{tabular}{@{}clrrrr@{}}
    \toprule
    \textbf{Run} & \textbf{LLM} & $\boldsymbol{N}$ & \textbf{Rec.} & \textbf{Rer.} & \textbf{LLM} \\
    \midrule
    1 & Qwen3.5-27B   & 309 & 42.07 & 56.31 & \textbf{81.23} \\
    2 & GPT-4o        & 308 & 42.53 & 56.82 & \textbf{81.82} \\
    4 & Qwen3.5-27B (FT joint) & 309 & 42.07 & 92.88 & \textbf{85.76} \\
    \bottomrule
  \multicolumn{6}{@{}l@{}}{\footnotesize Run~4: frozen \code{out\_bge\_rerank\_ft} $\rightarrow$ \code{cs\_llm\_result}, top-20; \code{Paper/data/run4\_joint\_ft\_measured.json}. Stock joint ref.\ 67.99\% ($N{=}303$) in Section~\ref{sec:exp-joint}.} \\
  \end{tabular}}
\end{table}

\begin{table*}[t]
  \caption{Issue-ID Hit@$K$ by pipeline stage. Recall: concatenated multi-turn user queries. LLM: last user turn only.}
  \label{tab:exp-hitk}
  \centering
  \footnotesize
  \setlength{\tabcolsep}{3pt}
  \renewcommand{\arraystretch}{1.05}
  \adjustbox{width=\textwidth,center}{%
  \begin{tabular}{@{}ll*{7}{r}@{}}
    \toprule
    Run & Stage & @1 & @3 & @5 & @10 & @15 & @20 & @50 \\
    \midrule
    1 & Hybrid recall & 42.07 & 61.49 & 74.43 & 84.47 & 88.35 & 90.29 & 96.76 \\
    1 & Rerank (top-50) & 56.31 & 75.80 & 84.14 & 91.59 & 94.50 & 95.15 & --- \\
    1 & LLM (Qwen3.5-27B) & \textbf{81.23} & --- & --- & --- & --- & --- & --- \\
    \midrule
    2 & Hybrid recall & 42.53 & 61.04 & 74.03 & 83.77 & 87.66 & 89.94 & 96.75 \\
    2 & Rerank (top-50) & 56.82 & 75.00 & 83.77 & 91.56 & 94.48 & 95.13 & --- \\
    2 & LLM (Azure GPT-4o) & \textbf{81.82} & --- & --- & --- & --- & --- & --- \\
    \bottomrule
  \end{tabular}}
\end{table*}

\paragraph{Stage-wise interpretation (Runs~1--2).}
The three-stage funnel in Figure~\ref{fig:performance_funnel} and Table~\ref{tab:exp-hitk} decomposes where error remains under a \emph{stock} reranker:
\begin{itemize}[leftmargin=1.2em]
  \item \textbf{Recall ($\approx$42\% Hit@1, $\approx$97\% Hit@50):} hybrid fusion finds the gold FAQ in the top-50 pool for most cases but rarely at rank~1---the system should not terminate at recall.
  \item \textbf{Rerank ($\approx$57\% Hit@1):} generic cross-encoder reranking helps (+15pp) yet $\approx$44\% of cases still have a wrong top-1 candidate.
  \item \textbf{Candidate election ($\approx$81\% Hit@1):} the decision-stage LLM reads candidate summaries and the latest user turn, recovering +25pp over rerank---primary offline evidence that the runtime should not stop at retrieve-then-rank when head ranking is weak.
\end{itemize}
Run~2 reproduces Run~1 with a different decision LLM; stable recall/rerank and near-identical LLM Hit@1 confirm that the funnel shape is a property of the \emph{pipeline}, not an artifact of a single vendor model.

\paragraph{Candidate window (top-10 vs.\ top-20).}
Run~1 reports decision-stage Hit@1 with \textbf{top-20} reranked candidates, matching the historical \code{workflow\_simple} archive and production eval defaults for this benchmark.
A controlled ablation on the \emph{same live stack} (May7 prompt, stock \code{bge-reranker-large}, Qwen3.5-27B; only \code{kb\_top\_k} changes) lowers the window to top-10 and yields \textbf{78.96\%} ($-$2.27pp; Table~\ref{tab:exp-topk}, rows~1--2).
Under weak head rerank (rerank@10 $=$ 89.97\% on frozen scores; Table~\ref{tab:rerank-biz}), exposing ranks 11--20 helps the election LLM recover cases where the gold FAQ sits just outside the top-10 cutoff.
Production \code{workflow\_simple} defaults to top-10 for latency; the ablation quantifies the accuracy--cost trade-off under stock rerank.

\begin{table}[t]
  \caption{Candidate-window sensitivity (Qwen3.5-27B decision, $N{=}309$). Stock rows: live \code{workflow\_simple} + May7 prompt. FT rows: frozen CUDA rerank scores + legacy LLM format (complementary protocol; Section~\ref{sec:exp-joint}).}
  \label{tab:exp-topk}
  \centering
  \footnotesize
  \setlength{\tabcolsep}{2pt}
  \adjustbox{width=\columnwidth,center}{%
  \begin{tabular}{@{}lccrr@{}}
    \toprule
    \textbf{Reranker} & \textbf{Path} & \textbf{Top-$K$} & \textbf{LLM @1} & $\boldsymbol{\Delta}$ \\
    \midrule
    stock bge & live & 20 & \textbf{81.23} & ref. \\
    stock bge & live & 10 & \textbf{78.96} & $-$2.27 \\
    \code{out\_bge\_rerank\_ft} & frozen CUDA & 10 & \textbf{85.76} & $+$4.53 vs.\ Run~1 \\
    \code{out\_bge\_rerank\_ft} & frozen CUDA & 20 & \textbf{85.76} & 0 vs.\ top-10 \\
    \bottomrule
  \multicolumn{5}{@{}l@{}}{\footnotesize Stock CSVs: \code{Paper/data/workflow\_simple\_...\_chain\_aligned\_20\_live.csv} and \code{...\_10.csv}. FT CSVs: \code{...\_ft\_rerank\_qwen\_top10\_legacy.csv}, \code{...\_top20\_legacy.csv}.} \\
  \end{tabular}}
\end{table}

\subsection{Reranker Fine-Tuning (Run~3)}
\label{sec:exp-rerank}

\noindent\textit{Purpose.} Run~3 responds to the bottleneck diagnosed in Runs~1--2: decision LLMs cannot fully compensate for weak head ranks when $\approx$44\% of top-1 rerank results are wrong. We therefore ask whether a reranker fine-tuned on in-domain FAQ pairs can close the $\approx$15-point gap between hybrid recall Hit@1 ($\approx$42\%) and off-the-shelf rerank Hit@1 ($\approx$57\%)---\emph{without} involving the LLM---and at what cost to public reranking benchmarks. After training, we evaluate each checkpoint on (a)~public reranking benchmarks and (b)~the 309-case business set under the recall-then-rerank protocol above.

\paragraph{General benchmark sanity check.}
We also checked public reranking benchmarks to avoid choosing a checkpoint that only memorizes the business set. The selected \code{out\_bge\_rerank\_ft} slightly lowers the average score relative to the public \code{bge-reranker-large} baseline, while the MMarco-mixed variant causes larger benchmark regression despite marginally higher business P@1. We therefore deploy the business fine-tune as the better production trade-off.

\paragraph{Business issue-ID precision.}
Table~\ref{tab:rerank-biz} reports precision at rank $K$ (P@$K$) on the 309 labeled cases after hybrid recall (top-50) and reranking---i.e., whether the gold \code{issue\_id} appears in the top-$K$ reranked list. The off-the-shelf reranker achieves P@1 57.28\%, consistent with Run~1 rerank Hit@1 (56.31\%). Fine-tuning on all 309 pairs lifts P@1 to \textbf{92.88\%} (+35.6pp) and saturates P@3--P@20 at $\geq$96.1\%. The 95/5 split variant reaches 89.32\% P@1; the MMarco-mixed variant peaks at 94.17\% P@1 but is not preferred given benchmark regression. Mean rerank latency remains $\approx$780\,ms (p95 $\approx$815\,ms), showing that domain adaptation does not inflate serving cost.

\begin{table*}[t]
  \caption{Business FAQ precision after recall (top-50) + rerank ($N{=}309$). Latency in ms per request.}
  \label{tab:rerank-biz}
  \centering
  \footnotesize
  \setlength{\tabcolsep}{3pt}
  \renewcommand{\arraystretch}{1.05}
  \adjustbox{width=\textwidth,center}{%
  \begin{tabular}{@{}l@{\hspace{2pt}}*{6}{c}@{\hspace{4pt}}cc@{}}
    \toprule
    \textbf{Reranker} & \textbf{P@1} & \textbf{P@3} & \textbf{P@5} & \textbf{P@10} & \textbf{P@15} & \textbf{P@20} & \textbf{mean} & \textbf{p95} \\
    \midrule
    bge-reranker-large & 57.28 & 75.73 & 83.82 & 89.97 & 95.15 & 95.15 & 780.0 & 815.3 \\
    \code{out\_bge\_rerank\_ft} & \textbf{92.88} & 96.12 & 96.76 & 96.76 & 96.76 & 96.76 & 779.1 & 814.8 \\
    \code{out\_ft\_split} & 89.32 & 95.47 & 96.76 & 96.76 & 96.76 & 96.76 & 778.8 & 814.0 \\
    \code{mixed\_marco\_r20} & 94.17 & 96.12 & 96.12 & 96.44 & 96.76 & 96.76 & 779.9 & 815.1 \\
    \bottomrule
  \end{tabular}}
\end{table*}

\paragraph{Implications (Run~3).}
A modest in-domain reranker fine-tune (+35.6pp P@1 on business data) exceeds the entire lift attributable to swapping Qwen for GPT-4o in Runs~1--2 ($<$1pp), while mean latency stays $\approx$780\,ms. We deploy \code{out\_bge\_rerank\_ft} in production: strong business head rank with acceptable benchmark drift. Mixing large MMarco corpora is avoided because it destroys in-domain ranking despite marginally higher P@1 on a handful of cases.

\subsection{Joint Rerank + LLM Pipeline (Run~4)}
\label{sec:exp-joint}

\noindent\textit{Purpose.} Runs~1--3 vary one stage at a time. Run~4 executes the \emph{full} offline pipeline in one pass---hybrid recall, cross-encoder rerank, then decision LLM---under the production fine-tuned checkpoint (Table~\ref{tab:canonical-metrics}).

\paragraph{Measured joint (fine-tuned rerank).}
We ran \code{run\_jsonl\_ft\_rerank\_llm\_issueid.py} with frozen CUDA scores from \code{out\_bge\_rerank\_ft}~\cite{xiao2024bge} and Qwen3.5-27B on $N{=}309$ labeled cases (May7 legacy prompt; concatenated recall query; last user turn for decision).
Table~\ref{tab:exp-joint} and \code{Paper/data/run4\_joint\_ft\_measured.json} report Hit@1 at each stage: recall 42.07\%, rerank 92.88\%, LLM \textbf{85.76\%} (+4.53pp over Run~1's 81.23\% anchor).
Marginal election lift shrinks to $\approx$+4.5pp once rerank head quality is strong---consistent with the ranker-first narrative---rather than the +25pp observed under stock rerank in Runs~1--2.

\paragraph{Candidate window under FT rerank.}
Table~\ref{tab:exp-topk} (rows~3--4) ablates the number of candidates passed to the decision LLM while holding frozen FT rerank scores fixed.
Both top-10 and top-20 yield \textbf{85.76\%} (265/309): rerank@10 already equals rerank@20 at 96.76\% (Table~\ref{tab:rerank-biz}), so ranks 11--20 add no new gold issues and the election LLM is \emph{insensitive} to the window width.
This justifies reporting Run~4 with top-20 for alignment with Run~1 while noting that production's top-10 default would not change FT joint accuracy on this benchmark.

\paragraph{Stock joint replication (reference).}
A separate one-pass stock pipeline via \code{paper\_faq\_pipeline\_eval.py} on $N{=}303$ CONTINUE-labeled cases logs LLM Hit@1 \textbf{67.99\%} (\code{Paper/data/run4\_joint\_measured\_v2.json})---below the Run~1 live archive (81.23\%) due to script filters, $N{=}303$, and harness differences (Section~\ref{sec:canonical-metrics}).
We retain it as reproducibility evidence that election still repairs wrong rerank heads (+15pp over rerank at $\approx$53\% head quality) but do not use it as the headline joint metric.
An earlier analytic projection ($\approx$97\% under optimistic repair transfer from Run~1 conditional rates) is superseded by the measured 85.76\% FT joint pass.

\begin{table}[t]
  \caption{Run~4 joint pipeline (Qwen3.5-27B decision). Row~1: measured FT joint; row~2: Run~1 stock anchor; row~3: stock joint script reference.}
  \label{tab:exp-joint}
  \centering
  \footnotesize
  \setlength{\tabcolsep}{3pt}
  \adjustbox{width=\columnwidth,center}{%
  \begin{tabular}{@{}lrrrr@{}}
    \toprule
    \textbf{Setting} & $\boldsymbol{N}$ & \textbf{Rec.} & \textbf{Rer.} & \textbf{LLM} \\
    \midrule
    FT rerank (meas.)     & 309 & 42.07 & 92.88 & \textbf{85.76} \\
    Run~1 stock (anchor)  & 309 & 42.07 & 56.31 & 81.23 \\
    Stock joint (ref.)    & 303 & 35.31 & 52.81 & 67.99 \\
    \bottomrule
  \end{tabular}}
\end{table}

\paragraph{Implications (Run~4).}
Measured FT joint confirms that domain rerank adaptation (Run~3) translates to end-to-end gains (+4.5pp over Run~1) without requiring a larger decision LLM.
Candidate-window insensitivity under FT rerank (Table~\ref{tab:exp-topk}) further supports deploying top-10 in production for latency once \code{out\_bge\_rerank\_ft} is active, while stock rerank deployments benefit from the wider window ($+$2.3pp for top-20 vs.\ top-10).

\subsection{Synthesis Across Runs}
\label{sec:exp-synthesis}
Table~\ref{tab:exp-synthesis} consolidates the experimental narrative. Runs~1--2 establish that \textbf{candidate election is necessary but not sufficient} when rerank head quality is poor: the system reaches $\approx$81\% only after the LLM repairs $\approx$44\% of wrong top-1 rerank outcomes. Run~3 shows that \textbf{improving evidence quality upstream} is the higher-leverage engineering investment on this dataset (+36pp at rerank vs.\ +25pp from LLM over stock rerank). Run~4 (Table~\ref{tab:exp-joint}) \emph{measures} FT joint at 85.8\% LLM Hit@1 ($+$4.5pp vs.\ Run~1); candidate-window ablation (Table~\ref{tab:exp-topk}) shows top-10 vs.\ top-20 matters under stock rerank ($-$2.3pp) but not under FT rerank (85.76\% for both). D2 and full DAG action accuracy remain unscored.

\begin{table}[t]
  \caption{Synthesis: what each experimental run establishes.}
  \label{tab:exp-synthesis}
  \centering
  \footnotesize
  \setlength{\tabcolsep}{3pt}
  \adjustbox{width=\columnwidth,center}{%
  \begin{tabular}{@{}>{\raggedright\arraybackslash}p{0.26\columnwidth}>{\raggedright\arraybackslash}p{0.70\columnwidth}@{}}
    \toprule
    \textbf{Evidence} & \textbf{Conclusion} \\
    \midrule
    Runs~1--2 funnel & Recall supplies a rich pool; stock rerank is weak at rank~1; LLM election adds +25pp. \\
    Run~1 vs.\ Run~2 & Backbone swap ($<$1pp) $\Rightarrow$ orchestration $>$ generator scale. \\
    Run~3 vs.\ stock rerank & Domain rerank FT adds +36pp P@1 at unchanged latency. \\
    Runs~1--2 + Run~3 & \textbf{Primary evidence:} D1 election + D3 rerank FT. \\
    Run~4 joint & FT joint measured 85.8\% Hit@1 (+4.5pp vs.\ Run~1); top-10/20 insensitive under FT rerank. \\
    Table~\ref{tab:exp-topk} & Stock: top-20 $+$2.3pp over top-10; FT: window width irrelevant once rerank@10 saturates. \\
    D2 (invariant) / full DAG & D2 not scored; MCP/action metrics planned (Section~\ref{sec:eval-scope}). \\
    \bottomrule
  \end{tabular}}
\end{table}

\subsection{Additional Evaluation (Planned)}
\label{sec:exp-planned}
We are extending evaluation along: (i)~\textbf{full runtime:} action accuracy with MCP and rule-as-evidence fusion; (ii)~\textbf{ablations and SLOs:} hard rule short-circuit vs.\ LLM fusion, dynamic routing, RRF-only vs.\ rerank, simulated channel timeouts, and latency p50/p95; (iii)~\textbf{live FT joint:} reproduce Run~4 under live HTTP \code{out\_bge\_rerank\_ft} (Run~4 today uses frozen CUDA scores).

\section{Discussion}
\subsection{What Is Novel Relative to Prior CS LLM Pipelines?}
\label{sec:discussion-novelty}
The paper's main novelty is not a new retriever, a larger chat model, or an unconstrained agent loop. It is the deployment pattern required for LLMs to take over customer-service decision chains without taking over production authority.
\textbf{Versus flat RAG and adaptive RAG}~\cite{lewis2020rag,asai2024selfrag,jeong2024adaptiverag,xu2024cskg}, our deployed object is not a generator prompt but a runtime harness: retrieval is one evidence source among FAQ, account tools, policy rules, user clarification, and human escalation.
\textbf{Versus single-loop tool agents}~\cite{yao2023react,singh2025agenticrag,qin2023toolllm}, D1 keeps execution authority in deterministic runtime code; the LLM can recommend bounded actions but cannot own an open-ended tool loop.
\textbf{Versus ReAct-style agents}, D2 turns ``ask the user'' into a governed action with slot state and counters; tool calls and user questions both pass through runtime constraints.
\textbf{Versus finite-state conversational routers}~\cite{sun2024dfarag}, we do not encode the whole conversation in a learned automaton; instead, we expose a fixed production DAG whose nodes assemble evidence, call tools, enforce rules, and decide when clarification or escalation is allowed.
\textbf{Versus rule-first bots}, D2 turns rules into auditable evidence candidates and forbids returning \code{rule\_answer} directly.
\textbf{Versus generic retrieve-then-rerank systems}, D3 treats ranking as one stage in a runtime readiness funnel that also includes clarification/action correctness and online outcomes.
The new evidence in this paper is that the deployed harness produces traceable stage-wise failures: Runs~1--2 show +25pp candidate-election lift under stock rerank, Run~3 shows +36pp from domain rerank adaptation at unchanged latency, and Run~4 measures FT joint at 85.76\% (+4.5pp vs.\ Run~1) with candidate-window sensitivity quantified in Table~\ref{tab:exp-topk}. D2 and full action metrics are supported by the runtime but require the planned rule/MCP/action labels in Section~\ref{sec:exp-planned}.

\subsection{Why Runtime Harnessing Matters}
Standard RAG baselines treat the LLM as a generator after fixed retrieval, or stop after cross-encoder reranking~\cite{nogueira2019bert}.
In a production takeover, however, the harder problem is not only what the LLM says but what the runtime allows it to see, ask, call, choose, escalate, and log.
Our Runs~1--2 show that even with $\approx$97\% Hit@50 recall, $\approx$57\% rerank Hit@1 leaves substantial error that candidate election can repair; Run~3 shows that upstream evidence quality is a higher-leverage intervention than swapping the chat backbone.
Under fine-tuned rerank, the marginal LLM gain shrinks to $\approx$+4.5pp (Run~4 measured joint in Section~\ref{sec:exp-joint}), reinforcing the ranker-first narrative for production prioritization.
This supports the harness view: model capability is expressed through runtime evidence quality and action boundaries, and trace attribution tells engineers whether to improve RAG, rerank, rules, clarification policy, or the model itself.

\subsection{What Makes the Patterns Reusable?}
D1--D3 are intentionally framed above our implementation. A different enterprise support team can reuse D1 by putting its LLM behind deterministic gates for identity, tool eligibility, slot sufficiency, escalation, timeout, and audit. It can reuse D2 by compiling every evidence channel---FAQ, account API, policy rule, transaction status, compliance flag---into typed candidates with provenance and allowed actions. It can reuse D3 by logging each turn so that failures map to update targets instead of disappearing into a generic ``bad answer'' bucket.

The trade-off is that the system is less flexible than a fully autonomous agent. The benefit is that the organization can decide which authority to expand after evidence improves. Our experiments illustrate this discipline: before giving the LLM more freedom, the trace funnel showed that rerank quality was the highest-leverage bottleneck, and domain rerank adaptation improved the evidence layer more than swapping Qwen for GPT-4o.

\section{Limitations and Future Work}
Experiments focus on FAQ issue-ID grounding on $N{\approx}309$ sessions from one financial CS domain---smaller than web-scale retrieval suites such as BEIR~\cite{thakur2021beir} but aligned with available labeled CS inventory; MCP tools, rule-derived evidence, clarification, and dialogue-action labels are deployed runtime behavior but not yet labeled in the offline set.
Run~3 trains and evaluates rerankers on the same 309 business cases (full-data fine-tuning), which may introduce optimistic bias in rerank P@1 relative to held-out production traffic---the 95/5 split variant (89.3\% P@1) partially mitigates this concern~\cite{xiao2024bge}.
Canonical headline metrics follow Runs~1--3 (Table~\ref{tab:canonical-metrics}); Run~4 reports measured FT joint accuracy (85.76\%) plus candidate-window sensitivity (Table~\ref{tab:exp-topk}); a stock joint script reference (67.99\%, $N{=}303$) is retained for reproducibility only.
For KDD Applied Data Science submission, the deployment table in Section~\ref{sec:deployment} should be filled with measured post-launch performance over a fixed traffic window.
We do not claim autonomous ``agent'' benchmark performance~\cite{singh2025agenticrag}; the contribution is a governed LLM runtime harness~\cite{langgraph2024,chung2023instructtods,zhou2026externalization} with explicit evaluation boundaries.
Generalization beyond financial-platform customer service remains to be tested.
Future work includes live FT joint reproduction, clarification/action/rule evidence metrics, cross-domain transfer, and formal analysis of candidate fusion under noisy rule triggers.

\section*{Ethics Statement}
Customer-service automation should include explicit escalation safeguards, policy-compliance checks, and logging for post-hoc auditability. Sensitive user attributes should be minimized in retrieval features and logs. Automated decisions that affect fund access or account security require human-in-the-loop escalation paths.

\bibliographystyle{ACM-Reference-Format}
\bibliography{refs}

\end{document}
